\documentclass[noauthor,noinstructornotes]{ximera}

\title{Week 7}
\author{Math 1193 Design Team}

\begin{document}
\begin{abstract}
    Activities for Week 7.
\end{abstract}
\maketitle

\begin{problem}
Find the following without a calculator. Write your answer as a 
mixed number (for example, \(6 \frac{3}{11}\)). Compare your method 
with your group.

\begin{enumerate}
    \item \(16 \div 3\)
    \item \(242 \div 9\)
    \item \(242 \div 19\)
\end{enumerate}
\end{problem}

\begin{problem}
In the equation \(35 \div 5 = 6 + \frac{4}{5}\), which number
is the divisor, which is the dividend, which is the quotient, and 
which is the remainder?
\end{problem}

\begin{problem}
Check that \(\frac{x^2 + x + 4}{x - 2} = x + 3 + \frac{10}{x - 2}\) by
rewriting the right-hand side using a common denominator. 

Identify the divisor, dividend, quotient, and remainder.

\end{problem}

\begin{problem}
\begin{enumerate}
    \item Check that \(x^2 - 1 = (x + 2)(x - 2) + 3\) by rewriting the right-hand
    side. You'll need to multiply and collect like terms.
    \item What is the result when plugging in \(x = 2\) to the right-hand side?
    \item Therefore, what is the result when plugging in \(x = 2\) to the left-hand
    side? You shouldn't have to do any extra work here, since the
    left- and right-hand sides are equal.
\end{enumerate}

\end{problem}

\begin{problem}
Rational functions are fractions where the numerator and
denominator are polynomials.

Is \(f(x) = \frac{2}{2 - x} - \frac{1}{x - 4}\) a rational function?
It's a \textbf{sum} of two fractions of polynomials, but can it be
written as just \textbf{one} fraction of polynomials?

Write \(f(x)\) as one fraction of polynomials.
\end{problem}

\begin{problem}
Let \(C\) be the function that gives the average cost per book 
\(C\left(x\right)\), in dollars, when using an online store to print
\(x\) copies of a self-published paperback book. Here is a graph of 
\(c\left(x\right)=\frac{120+4x}{x}\). 

\begin{center}
    \includegraphics[width=20cm]{bookselling.png}
\end{center}

\begin{enumerate}
    \item What is the approximate cost per book when 50 books are 
    printed? 100 books?
    \item The author plans to charge \$8 per book. About how many
    should be printed to make a profit?
    \item What is the value of \(C(x)\) when \(x = \frac{1}{2}\)? how
    does this relate to the context?
    \item What does the end behavior of the function \(C\) say about
    the context. Write down your answer, then share with a partner.
\end{enumerate}

\begin{instructorNotes}
For the third part, invite the discussion of the ``half-book''. Does
selling a half-book make sense? What about a book that's half as long?
\end{instructorNotes}
\end{problem}

\begin{problem}
For each of the following functions, solve the equation \(f(x) = 0\)
and evaluate \(f(0)\). What do these answers tell you about the graph 
of \(f\)?

\begin{enumerate}
    \item \(f(x) = \frac{x^2 - 1}{x + 2}\)
    \item \(f(x) = \frac{x + 1}{x^2 - 1}\)
    \item \(f(x) = \frac{1}{x - 3} + \frac{2}{x + 1}\)
    \item \(f(x) = \frac{2x + 1}{x^ - x}\)
\end{enumerate}

\begin{instructorNotes}
This is meant to see if students recognize the idea of intercepts and how
they connect to the formula.
\end{instructorNotes}
\end{problem}

\end{document}
